\subsection{Semisimple Rings}

THEOREM 1 (Wedderburn). --- Let A be a ring. The following properties are equivalent:

\begin{enumerate}
\def\labelenumi{(\roman{enumi})}
\item
  The A-module \(\mathrm{A}_s\) is semisimple.
\item
  For every left ideal \(\mathfrak{a}\) of A, there exists a left ideal \(\mathfrak{b}\) of A such that \(\mathrm{A}_s\) is the direct sum of \(\mathfrak{a}\) and \(\mathfrak{b}\) .
\item
  The ring A is left Artinian, and the (A,A)-bimodule \(_{s}A_{d}\) is semi-simple.
\item
  The ring A is isomorphic to the product of a finite family of simple rings.
\item
  There exist an integer \(s \geqslant 0\) , fields \(\mathrm{D}_1, \ldots, \mathrm{D}_s\) , and integers \(r_1 \geqslant 1, \ldots, r_s \geqslant 1\) such that the ring A is isomorphic to the product of the matrix rings \(\mathbf{M}_{r_i}(\mathrm{D}_i)\) .
\item
  The ring A is left Artinian, and there exists a faithful and semisimple left A-module.
\end{enumerate}

The equivalence of (i) and (ii) follows from Corollary 2 of VIII, p.~56, and that of (iv) and (v) follows from Theorem 1 of VIII, p.~120.

The A-module \(A_{s}\) is finitely generated. If it is semisimple, then it is left Artinian (VIII, p.~71, Proposition 10). Since the A-module \(A_{s}\) is faithful, this proves that (i) implies (vi). Conversely, suppose that property (vi) holds; let M be a faithful semisimple A-module. There exists an integer \(m \geqslant 1\) such that \(A_{s}\) is isomorphic to a submodule of \(M^{m}\) (VIII, p.~50, Proposition 5, a)); since M is semisimple, the same holds for \(A_{s}\) . We have therefore proved the equivalence of (i) and (vi).

Let us prove that (i) implies (iii). Suppose that the ring A has property (i); we have already observed that A is then left Artinian. Now, the endomorphisms of the left A-module \(A_{s}\) are the right multiplications by the elements of A. That the (A, A)-bimodule \(_{s}A_{d}\) is semisimple then follows from Proposition 6 of VIII, p.~86.

Let us show that (iii) implies (iv). Suppose that the (A, A)-bimodule \(_{s}A_{d}\) is semisimple. It is finitely generated, so there exists a finite family \((\mathfrak{a}_{i})_{i\in I}\) of simple (A, A)-sub-bimodules with direct sum \(_{s}A_{d}\) . In other words, the \(a_{i}\) are nonzero two-sided ideals of A, the additive group of A is the direct sum of the \(a_{i}\) , and for every \(i\in I\) , every two-sided ideal of A contained in \(a_{i}\) is equal to 0 or \(a_{i}\) . Set \(b_{i}=\sum_{j\neq i}a_{j}\) for every \(i\in I\) ; this is a two-sided ideal of A. The mapping \(a\mapsto(a+\mathfrak{b}_{i})_{i\in I}\) is an isomorphism from the ring A to the product of the rings \(A/b_{i}\) . The (A, A)-bimodules \(a_{i}\) and \(A/b_{i}\) are isomorphic, so every two-sided ideal of \(A/b_{i}\) is equal to 0 or \(A/b_{i}\) . If the ring A is left Artinian, then the same holds for the rings \(A/b_{i}\) , which are therefore simple (VIII, p.~120, Definition 1).

Finally, let us prove that (iv) implies (i). Suppose that A is the product of a finite family \((\mathrm{A}_{i})_{i\in\mathrm{I}}\) of simple rings. Denote by \(\pi_{i}\) the projection with index i from A to \(A_{i}\) and by \(M_{i}\) the A-module with underlying additive group \(A_{i}\) and law of action \((a,x)\mapsto\pi_{i}(a)x\) . Since the ring \(A_{i}\) is simple, the \(A_{i}\) -module \((\mathrm{A}_{i})_{s}\) is semisimple, so the A-module \(M_{i}\) is semisimple. Since the A-module \(A_{s}\) is nothing but the product \(\prod_{i\in I}M_{i}\) , it is semisimple.

DEFINITION 1. --- We say that a ring A is semisimple if it has the equivalent properties (i) through (vi) of Theorem 1. An algebra A over a commutative ring k is a semisimple algebra if the ring underlying A is semisimple.

PROPOSITION 1. --- Let A be a semisimple ring. There exists a finite family \((\mathfrak{m}_{i})_{i\in I}\) of minimal left ideals of A such that \(A_{s}=\oplus_{i\in I}m_{i}\) . If \((\mathfrak{m}_{i})_{i\in I}\) is such a family, then every simple A-module is isomorphic to one of the \(m_{i}\) . The set of classes of simple A-modules is finite.

The first assertion follows from the fact that the A-module \(A_{s}\) is semisimple and finitely generated. Every simple module is isomorphic to a quotient of \(A_{s}\) (VIII, p.~46, Proposition 1). The second assertion then follows from Corollary 3 of VIII, p.~56, and the third assertion is an immediate consequence of the second.

Example. --- Let G be a finite group and K a commutative field. *We will see further on (VIII, p.~401, Corollary 1) that the algebra K{[}G{]} of the group G over the field K is a semisimple ring if and only if the characteristic exponent of K is prime to the order of G.*

Remarks. --- 1) Let K be a commutative field, and let A be a semisimple algebra over K. Then there exist K-algebras \(D_{1},\ldots,D_{s}\) that are fields and integers \(r_{1}\geqslant1,\ldots,r_{s}\geqslant1\) such that the K-algebra A is isomorphic to the product \(\prod_{i=1}^{s}\mathbf{M}_{r_{i}}(\mathrm{D}_{i})\) .

\begin{enumerate}
\def\labelenumi{\arabic{enumi})}
\setcounter{enumi}{1}
\tightlist
\item
  Let K be an algebraically closed field, and let A be an algebra of finite degree over K. By Remark 4 of VIII, p.~122, the algebra A is semisimple if and only if there exist integers \(n_{1} \geqslant 1, \ldots, n_{r} \geqslant 1\) such that A is isomorphic to the algebra \(\mathbf{M}_{n_{1}}(\mathrm{K}) \times \cdots \times \mathbf{M}_{n_{r}}(\mathrm{K})\) .
\end{enumerate}

PROPOSITION 2. --- a) The center of a semisimple ring is semisimple.

\begin{enumerate}
\def\labelenumi{\alph{enumi})}
\setcounter{enumi}{1}
\item
  The opposite ring of a semisimple ring is semisimple.
\item
  The quotient of a semisimple ring by a two-sided ideal is a semisimple ring.
\item
  The product of a finite family of semisimple rings is a semisimple ring.
\end{enumerate}

Let A be a semisimple ring. It is isomorphic to the product of a finite family \((\mathrm{A}_{i})_{i\in\mathrm{I}}\) of simple rings. The center of A is isomorphic to the product of the centers of the \(A_{i}\) , and \(A^{o}\) is isomorphic to the product ring of the \(A_{i}^{o}\) . Assertions a) and b) therefore follow from Corollary 1 of VIII, p.~121.

Let \(\mathfrak{a}\) be a two-sided ideal of A. The A-module \(A_s / \mathfrak{a}\) , quotient of the semisimple A-module \(A_s\) , is semisimple. The \((A / \mathfrak{a})\) -module \(A_s / \mathfrak{a}\) is therefore semisimple; assertion c) follows.

A ring is semisimple if and only if it is isomorphic to the product of a finite family of simple rings; assertion d) follows.

PROPOSITION 3. --- Let A be a commutative ring. The following properties are equivalent:

\begin{enumerate}
\def\labelenumi{(\roman{enumi})}
\item
  The ring A is semisimple.
\item
  The ring A is Artinian and reduced (V, §6, No.~7, p.~34).
\item
  The ring A is isomorphic to the product of a finite family of commutative fields.
\end{enumerate}

Commutative simple rings are commutative fields (VIII, p.~120, Remark 1). So (i) is equivalent to (iii).

It is clear that (iii) implies (ii). Conversely, suppose that the ring A is Artinian and reduced. The intersection of the set of prime ideals of A consists of the nilpotent elements of A (V, §15, No.~1, p.~118, Proposition 2), hence is reduced to 0 because A is reduced. By VIII, p.~2, there then exist distinct prime ideals \(\mathfrak{p}_1, \ldots, \mathfrak{p}_r\) of A such that we have \(\mathfrak{p}_1 \cap \cdots \cap \mathfrak{p}_r = 0\) . By the corollary of VIII, p.~8, each of the prime ideals \(\mathfrak{p}_i\) of the Artinian ring A is maximal; we therefore have \(\mathfrak{p}_i + \mathfrak{p}_j = A\) whenever \(i\) and \(j\) are distinct. By Proposition 9 of I, §8, No.~11, p.~110, the canonical homomorphism from A to the ring \(\prod_{i=1}^{r} (A / \mathfrak{p}_i)\) is an isomorphism. The ring \(A / \mathfrak{p}_i\) is a field for every \(i\) , and (ii) therefore implies (iii).

A commutative algebra of finite degree over a field is an Artinian commutative ring. Proposition 3 therefore generalizes Proposition 5 of V, §6, No.~7, p.~34.

